% Part: first-order-logic
% Chapter: syntax-and-semantics
% Section: covered-structures

\documentclass[../../../include/open-logic-section]{subfiles}

\begin{document}

\olfileid{fol}{syn}{cov}

\olsection{پہلے درجے دیاں زباناں لئی ڈھکیاں ہوئیاں \printtoken{P}{structure}}

\begin{explain}
یاد رکھو، اک ٹرم \emph{بند} اے جے اوہدے وچ کوئی !!{variable}s نہیں۔
\end{explain}

\begin{defn}[بند ٹرماں دی !!^{value}]
جے $t$، زبان~$\Lang L$ دا بند ٹرم اے تے $\Struct M$،
~$\Lang L$ لئی !!{structure} اے، تاں اوہدی
\emph{!!{value}}~$\Value{t}{M}$ ایویں تعریف ہوندی اے:
\begin{enumerate}
\item جے $t$ بس !!{constant}~$c$ اے، تاں $\Value{c}{M} = \Assign{c}{M}$۔
\item جے $t$ دی صورت $\Atom{f}{t_1, \ldots, t_n}$ اے، تاں
  \[
  \Value{t}{M} = \Assign{f}{M}(\Value{t_1}{M}, \ldots,
  \Value{t_n}{M}).
  \]
\end{enumerate}
\end{defn}

\begin{defn}[ڈھکی ہوئی !!{structure}]
!!{structure} \emph{ڈھکی ہوئی} اے جے دائرے دا ہر عنصر
کسے بند ٹرم دی !!{value} ہووے۔
\end{defn}

\begin{ex}
فرض کرو~$\Lang L$ ایسی زبان اے جیہدے وچ !!{constant}s
$\Obj{zero}$، $\Obj{one}$، $\Obj{two}$، \dots، دو-مقامی
!!{predicate}~$<$، تے دو-مقامی !!{function}s $+$ تے $\times$
نیں۔ فیر !!a{structure}~$\Struct
M$،~$\Lang L$ لئی، اوہ اے جیہدا دائرہ $\Domain M = \{0, 1, 2, \ldots
\}$ اے، تے تعبیراں $\Assign{\Obj{zero}}{M} = 0$،
$\Assign{\Obj{one}}{M} = 1$، $\Assign{\Obj{two}}{M} = 2$،
تے ایس طرح اگے نیں۔ دو-مقامی تعلق دی علامت $<$ لئی سیٹ
$\Assign{<}{M}$ اوہناں ساریاں جوڑیاں دا سیٹ اے
$\tuple{c_1, c_2} \in \Domain{M}^2$ جنہاں وچ
$c_1$،~$c_2$ توں گھٹ اے: مثال لئی، $\tuple{1, 3} \in
\Assign{<}{M}$ پر $\tuple{2, 2} \notin \Assign{<}{M}$۔
دو-مقامی !!{function} $+$ لئی، $\Assign{+}{M}$ نوں معمول
مطابق تعریف کرو---مثلاً، $\Assign{+}{M}(2,3)$ دا حاصل~$5$
اے؛ دو-مقامی !!{function}~$\times$ لئی وی ایویں۔ ایس لئی
$\Obj{four}$ دی !!{value} بس~$4$ اے، تے
$\times(\Obj{two},
+(\Obj{three},\Obj{zero}))$ دی !!{value} (یا درمیان-علامتی
لکھت وچ، $\Obj{two} \times
(\Obj{three} + \Obj{zero})$) ایہ اے:
% PNB-SEM-002: منجمد ماخذ وچ پہلی مساوات دی علامت دو واری اے؛ ایتھے صرف اگلی سطر دی اک علامت رکھی اے۔
\begin{multline*}
\Value{\times(\Obj{two}, +(\Obj{three},\Obj{zero}))}{M} \\
\begin{aligned}
& =\Assign{\times}{M}(\Value{\Obj{two}}{M}, \Value{+(\Obj{three}, \Obj{zero})}{M})\\
& = \Assign{\times}{M}(\Value{\Obj{two}}{M}, \Assign{+}{M}(\Value{\Obj{three}}{M},
\Value{\Obj{zero}}{M})) \\
& = \Assign{\times}{M}(\Assign{\Obj{two}}{M}, \Assign{+}{M}(\Assign{\Obj{three}}{M},
\Assign{\Obj{zero}}{M})) \\
& = \Assign{\times}{M}(2, \Assign{+}{M}(3, 0)) \\
& = \Assign{\times}{M}(2, 3) \\
& = 6
\end{aligned}
\end{multline*}
\end{ex}

\begin{prob}
کی حساب دا معیاری نمونہ $\Struct N$ ڈھکیا ہویا اے؟
اپنا جواب سمجھاؤ۔
\end{prob}

\end{document}
